% ============================================================================
% (c) 2025 Mohamed EL AFRIT -- IPSSI
% Licence CC BY-NC-SA 4.0
% https://creativecommons.org/licenses/by-nc-sa/4.0/deed.fr
%
% FICHE D'EXERCICES (enonces) -- Jour 2 : Optimisation
% ============================================================================
% ============================================================================
% © 2025 Mohamed EL AFRIT — IPSSI
% Ce contenu est distribué sous licence Creative Commons BY-NC-SA 4.0
% https://creativecommons.org/licenses/by-nc-sa/4.0/deed.fr
% ============================================================================
%
% TEMPLATE COMMUN — Mathématiques appliquées à l'Intelligence Artificielle
% Ce fichier est le préambule partagé de TOUS les documents LaTeX du projet.
%
% Utilisation :
%   % ============================================================================
% © 2025 Mohamed EL AFRIT — IPSSI
% Ce contenu est distribué sous licence Creative Commons BY-NC-SA 4.0
% https://creativecommons.org/licenses/by-nc-sa/4.0/deed.fr
% ============================================================================
%
% TEMPLATE COMMUN — Mathématiques appliquées à l'Intelligence Artificielle
% Ce fichier est le préambule partagé de TOUS les documents LaTeX du projet.
%
% Utilisation :
%   % ============================================================================
% © 2025 Mohamed EL AFRIT — IPSSI
% Ce contenu est distribué sous licence Creative Commons BY-NC-SA 4.0
% https://creativecommons.org/licenses/by-nc-sa/4.0/deed.fr
% ============================================================================
%
% TEMPLATE COMMUN — Mathématiques appliquées à l'Intelligence Artificielle
% Ce fichier est le préambule partagé de TOUS les documents LaTeX du projet.
%
% Utilisation :
%   \input{template_ipssi}          % Importer ce préambule
%   \jourNumero{1}                  % Définir le numéro du jour (1, 2, 3 ou 4)
%   \jourTheme{Données et représentation mathématique}
%   \begin{document}
%   \pagegarde{Fiche de synthèse}{1}{Données et représentation mathématique}
%   ...
%   \end{document}
% ============================================================================

% --- Classe de document ---
\documentclass[11pt,a4paper]{article}

% ============================================================================
% ENCODAGE ET LANGUE
% ============================================================================
\usepackage[utf8]{inputenc}
\usepackage[T1]{fontenc}
\usepackage[french]{babel}

% ============================================================================
% MATHÉMATIQUES
% ============================================================================
\usepackage{amsmath,amssymb,amsfonts,mathtools}

% Commandes mathématiques personnalisées (conventions du cours)
\newcommand{\vx}{\mathbf{x}}          % vecteur x
\newcommand{\vw}{\mathbf{w}}          % vecteur w (poids)
\newcommand{\vy}{\mathbf{y}}          % vecteur y
\newcommand{\mX}{\mathbf{X}}          % matrice X
\newcommand{\mW}{\mathbf{W}}          % matrice W
\newcommand{\yhat}{\hat{y}}           % prédiction
\newcommand{\pscal}[2]{\langle #1, #2 \rangle} % produit scalaire
\newcommand{\gradJ}{\nabla J(\boldsymbol{\theta})} % gradient de J

% ============================================================================
% MISE EN PAGE
% ============================================================================
\usepackage[margin=2.5cm]{geometry}
\usepackage{setspace}
\setstretch{1.15}                      % interligne légèrement aéré

% ============================================================================
% COULEURS
% ============================================================================
\usepackage[dvipsnames,svgnames,x11names]{xcolor}

% Palette principale du projet
\definecolor{ipssiBleu}{HTML}{3B82F6}       % Définitions
\definecolor{ipssiVert}{HTML}{22C55E}       % À retenir
\definecolor{ipssiOrange}{HTML}{F97316}     % Attention / Piège
\definecolor{ipssiViolet}{HTML}{A855F7}     % Intuition
\definecolor{ipssiCyan}{HTML}{06B6D4}       % Exemples
\definecolor{ipssiGrisClair}{HTML}{F3F4F6}  % Fond code
\definecolor{ipssiGrisFonce}{HTML}{1F2937}  % Texte code
\definecolor{ipssiFondPage}{HTML}{FAFAFA}   % Fond léger page de garde
\definecolor{ipssiRouge}{HTML}{E74C3C}      % Classe 0 (salaire ≤ 45k€)
\definecolor{ipssiVertData}{HTML}{2ECC71}   % Classe 1 (salaire > 45k€)
\definecolor{ipssiBleuDroite}{HTML}{3498DB} % Droite de régression

% ============================================================================
% ENCADRÉS PÉDAGOGIQUES (tcolorbox)
% ============================================================================
\usepackage[most]{tcolorbox}

% Style de base partagé
\tcbset{
    ipssibase/.style={
        enhanced,
        breakable,
        left=4mm,
        right=4mm,
        top=3mm,
        bottom=3mm,
        boxrule=0pt,
        frame hidden,
        borderline west={3pt}{0pt}{#1},
        colback=#1!5,
        colframe=#1,
        fonttitle=\bfseries\sffamily,
        coltitle=#1!80!black,
        attach boxed title to top left={yshift=-2mm, xshift=4mm},
        boxed title style={
            colback=#1!15,
            colframe=#1,
            boxrule=0.5pt,
            arc=2pt,
        },
        arc=2pt,
        outer arc=2pt,
    }
}

% Icônes TikZ pour les encadrés (compatible pdflatex, pas d'emoji Unicode)
\newcommand{\icoDefinition}{%
    \tikz[baseline=-0.3em]{\draw[ipssiBleu, thick] (0,0) -- (0.3,0) -- (0.3,0.35) -- (0,0.35) -- cycle;
    \draw[ipssiBleu, thick] (0.05,0.1) -- (0.25,0.1); \draw[ipssiBleu, thick] (0.05,0.2) -- (0.2,0.2);}%
}
\newcommand{\icoRetenir}{%
    \tikz[baseline=-0.2em]{\draw[ipssiVert, thick, fill=ipssiVert!20] (0.15,0) -- (0.3,0.3) -- (0,0.3) -- cycle;}%
}
\newcommand{\icoAttention}{%
    \tikz[baseline=-0.2em]{\draw[ipssiOrange, thick] (0.15,0.35) -- (0,0) -- (0.3,0) -- cycle;
    \node[ipssiOrange, font=\tiny\bfseries] at (0.15,0.12) {!};}%
}
\newcommand{\icoIntuition}{%
    \tikz[baseline=-0.2em]{\draw[ipssiViolet, thick, fill=ipssiViolet!15]
    (0.15,0.35) -- (0.08,0.2) -- (0.08,0.15) -- (0.22,0.15) -- (0.22,0.2) -- cycle;
    \draw[ipssiViolet, thick] (0.11,0.12) -- (0.11,0.05); \draw[ipssiViolet, thick] (0.19,0.12) -- (0.19,0.05);}%
}
\newcommand{\icoExemple}{%
    \tikz[baseline=-0.2em]{\draw[ipssiCyan, thick] (0,0) rectangle (0.3,0.3);
    \draw[ipssiCyan, thick] (0,0.15) -- (0.3,0.15); \draw[ipssiCyan, thick] (0.15,0) -- (0.15,0.3);}%
}
\newcommand{\icoPython}{%
    \tikz[baseline=-0.2em]{\node[font=\small\bfseries\ttfamily, ipssiVert!70!black] at (0,0) {\textgreater\textgreater};}%
}

% Définition (bleu)
\newtcolorbox{definitionbox}[1][D\'{e}finition]{
    ipssibase=ipssiBleu,
    title={\icoDefinition\; #1},
}

% À retenir (vert)
\newtcolorbox{retenirbox}[1][\`{A} retenir]{
    ipssibase=ipssiVert,
    title={\icoRetenir\; #1},
}

% Attention / Piège (orange)
\newtcolorbox{attentionbox}[1][Attention]{
    ipssibase=ipssiOrange,
    title={\icoAttention\; #1},
}

% Intuition (violet)
\newtcolorbox{intuitionbox}[1][Intuition]{
    ipssibase=ipssiViolet,
    title={\icoIntuition\; #1},
}

% Exemple (cyan)
\newtcolorbox{exemplebox}[1][Exemple]{
    ipssibase=ipssiCyan,
    title={\icoExemple\; #1},
}

% Code Python (style terminal)
\newtcolorbox{pythonbox}[1][Code Python]{
    enhanced,
    breakable,
    left=4mm, right=4mm, top=3mm, bottom=3mm,
    boxrule=0.5pt,
    colback=ipssiGrisClair,
    colframe=ipssiGrisFonce!40,
    fonttitle=\bfseries\ttfamily\small,
    coltitle=ipssiGrisFonce,
    title={\icoPython\; #1},
    attach boxed title to top left={yshift=-2mm, xshift=4mm},
    boxed title style={colback=ipssiGrisClair, colframe=ipssiGrisFonce!40, boxrule=0.5pt, arc=2pt},
    arc=2pt, outer arc=2pt,
}

% ============================================================================
% CODE PYTHON (listings)
% ============================================================================
\usepackage{listings}

\lstdefinestyle{python}{
    language=Python,
    basicstyle=\ttfamily\small,
    keywordstyle=\color{ipssiBleu}\bfseries,
    stringstyle=\color{ipssiVert!70!black},
    commentstyle=\color{gray}\itshape,
    numberstyle=\tiny\color{gray},
    numbers=left,
    numbersep=8pt,
    backgroundcolor=\color{ipssiGrisClair},
    frame=none,
    rulecolor=\color{ipssiGrisFonce!20},
    breaklines=true,
    breakatwhitespace=true,
    tabsize=4,
    showstringspaces=false,
    captionpos=b,
    aboveskip=8pt,
    belowskip=8pt,
    xleftmargin=10pt,
    framexleftmargin=10pt,
    morekeywords={as, with, True, False, None, self, np, plt, import, from},
    literate=
        {é}{{\'e}}1 {è}{{\`e}}1 {ê}{{\^e}}1 {ë}{{\"e}}1
        {à}{{\`a}}1 {â}{{\^a}}1
        {ù}{{\`u}}1 {û}{{\^u}}1
        {ô}{{\^o}}1
        {î}{{\^i}}1 {ï}{{\"i}}1
        {ç}{{\c{c}}}1
        {É}{{\'E}}1 {È}{{\`E}}1
        {→}{{$\rightarrow$}}1
        {≈}{{$\approx$}}1
        {≤}{{$\leq$}}1
        {≥}{{$\geq$}}1
        {★}{{$\bigstar$}}1
        {ŷ}{{$\hat{y}$}}1
        {·}{{$\cdot$}}1,
}

\lstset{style=python}

% Style pour le code inline
\newcommand{\pyinline}[1]{\lstinline[style=python,basicstyle=\ttfamily\small]{#1}}

% ============================================================================
% GRAPHIQUES
% ============================================================================
\usepackage{tikz}
\usepackage{pgfplots}
\pgfplotsset{compat=1.18}
\usetikzlibrary{arrows.meta, calc, positioning, shapes.geometric, decorations.pathreplacing}

% ============================================================================
% TABLEAUX
% ============================================================================
\usepackage{booktabs}
\usepackage{tabularx}
\usepackage{multirow}

% ============================================================================
% LISTES
% ============================================================================
\usepackage{enumitem}
\setlist[itemize]{itemsep=2pt, parsep=0pt, topsep=4pt}
\setlist[enumerate]{itemsep=2pt, parsep=0pt, topsep=4pt}

% ============================================================================
% HYPERLIENS
% ============================================================================
\usepackage[
    colorlinks=true,
    linkcolor=ipssiBleu!80!black,
    urlcolor=ipssiBleu!80!black,
    citecolor=ipssiVert!80!black,
    bookmarks=true,
    bookmarksopen=true,
    pdfauthor={Mohamed EL AFRIT},
    pdftitle={Mathématiques appliquées à l'IA — IPSSI},
    pdfsubject={Formation Bachelor 2 — IPSSI 2025-2026},
]{hyperref}

% ============================================================================
% IMAGES
% ============================================================================
\usepackage{graphicx}

% ============================================================================
% VARIABLES PARAMÉTRABLES (jour et thème)
% ============================================================================
\newcommand{\leJour}{X}
\newcommand{\leTheme}{Thème}
\newcommand{\jourNumero}[1]{\renewcommand{\leJour}{#1}}
\newcommand{\jourTheme}[1]{\renewcommand{\leTheme}{#1}}

% ============================================================================
% EN-TÊTE ET PIED DE PAGE (fancyhdr)
% ============================================================================
\usepackage{fancyhdr}

\pagestyle{fancy}
\fancyhf{}

% En-tête
\fancyhead[L]{\small\sffamily\textcolor{gray}{IPSSI — Maths appliquées à l'IA}}
\fancyhead[R]{\small\sffamily\textcolor{gray}{Jour \leJour\ — \leTheme}}
\renewcommand{\headrulewidth}{0.4pt}

% Pied de page
\fancyfoot[L]{\footnotesize\textcolor{gray}{Mohamed EL AFRIT\enspace\textbullet\enspace\url{www.mohamedelafrit.com}\enspace\textbullet\enspace CC BY-NC-SA 4.0}}
\fancyfoot[C]{\footnotesize\textcolor{gray}{\thepage}}
\fancyfoot[R]{\footnotesize\textcolor{gray}{2025–2026}}
\renewcommand{\footrulewidth}{0.2pt}

% Style pour la page de garde (pas d'en-tête)
\fancypagestyle{pagegarde}{
    \fancyhf{}
    \renewcommand{\headrulewidth}{0pt}
    \fancyfoot[C]{\footnotesize\textcolor{gray}{\thepage}}
    \renewcommand{\footrulewidth}{0pt}
}

% ============================================================================
% PAGE DE GARDE — \pagegarde{titre}{jour}{theme}
% ============================================================================
\newcommand{\pagegarde}[3]{%
    \thispagestyle{pagegarde}
    \begin{center}
        \vspace*{1.5cm}

        % Logo / Bandeau institution
        {\Large\sffamily\bfseries\textcolor{ipssiBleu}{IPSSI}}\par
        \vspace{2pt}
        {\normalsize\sffamily 2\textsuperscript{e} année Bachelor Informatique}\par

        \vspace{0.8cm}
        {\color{ipssiBleu}\rule{0.6\textwidth}{1.5pt}}\par
        \vspace{0.8cm}

        % Titre du module
        {\LARGE\sffamily\bfseries Mathématiques appliquées\\[4pt] à l'Intelligence Artificielle}\par

        \vspace{1cm}

        % Titre du document
        \begin{tcolorbox}[
            enhanced,
            colback=ipssiBleu!5,
            colframe=ipssiBleu!50,
            boxrule=1pt,
            arc=4pt,
            width=0.85\textwidth,
            halign=center,
            top=10pt, bottom=10pt,
        ]
            {\huge\sffamily\bfseries\textcolor{ipssiBleu!80!black}{#1}}\par
            \vspace{6pt}
            {\Large\sffamily Jour #2 — #3}
        \end{tcolorbox}

        \vspace{1.5cm}

        % Informations auteur
        {\large\sffamily
            \textbf{Auteur :} Mohamed EL AFRIT\\[4pt]
            \url{www.mohamedelafrit.com}\\[8pt]
            \textbf{Année académique :} 2025–2026
        }\par

        \vspace{1.5cm}

        % Licence Creative Commons
        {\color{ipssiBleu}\rule{0.4\textwidth}{0.5pt}}\par
        \vspace{8pt}
        {\small\sffamily
            \textcopyright\ 2025 Mohamed EL AFRIT — IPSSI\\[2pt]
            Ce contenu est distribué sous licence\\[2pt]
            \textbf{Creative Commons BY-NC-SA 4.0}\\[2pt]
            \url{https://creativecommons.org/licenses/by-nc-sa/4.0/deed.fr}
        }\par
    \end{center}
    \newpage
}

% ============================================================================
% NIVEAUX D'EXERCICES
% ============================================================================

% Étoile compatible pdflatex
\newcommand{\starfull}[1]{\textcolor{#1}{$\bigstar$}}
\newcommand{\starempty}{\textcolor{gray!40}{$\bigstar$}}

% Fondamental (1 étoile)
\newcommand{\exofondamental}{%
    \starfull{ipssiVert}\starempty\starempty\enspace%
    {\small\sffamily\textcolor{ipssiVert}{\textbf{Fondamental}}}%
}

% Intermédiaire (2 étoiles)
\newcommand{\exointermediaire}{%
    \starfull{ipssiOrange}\starfull{ipssiOrange}\starempty\enspace%
    {\small\sffamily\textcolor{ipssiOrange}{\textbf{Interm\'{e}diaire}}}%
}

% Avancé (3 étoiles)
\newcommand{\exoavance}{%
    \starfull{ipssiRouge}\starfull{ipssiRouge}\starfull{ipssiRouge}\enspace%
    {\small\sffamily\textcolor{ipssiRouge}{\textbf{Avanc\'{e}}}}%
}

% Commande d'exercice numéroté
\newcounter{exocount}[section]
\newcommand{\exercice}[1]{%
    \stepcounter{exocount}%
    \vspace{8pt}%
    \noindent\textbf{\sffamily Exercice \theexocount}\enspace — \enspace #1\par
    \vspace{4pt}%
}

% ============================================================================
% COMMANDES UTILITAIRES
% ============================================================================

% Séparateur visuel entre sections
\newcommand{\separateur}{%
    \vspace{6pt}
    \begin{center}
        {\color{ipssiBleu!30}\rule{0.3\textwidth}{0.5pt}}
    \end{center}
    \vspace{6pt}
}

% Mot-clé en gras coloré
\newcommand{\motcle}[1]{\textbf{\textcolor{ipssiBleu!80!black}{#1}}}

% Formule encadrée importante
\newcommand{\formule}[1]{%
    \begin{center}
        \fcolorbox{ipssiBleu!50}{ipssiBleu!5}{%
            \quad$\displaystyle #1$\quad%
        }
    \end{center}
}

% ============================================================================
% FIN DU PRÉAMBULE
% ============================================================================
          % Importer ce préambule
%   \jourNumero{1}                  % Définir le numéro du jour (1, 2, 3 ou 4)
%   \jourTheme{Données et représentation mathématique}
%   \begin{document}
%   \pagegarde{Fiche de synthèse}{1}{Données et représentation mathématique}
%   ...
%   \end{document}
% ============================================================================

% --- Classe de document ---
\documentclass[11pt,a4paper]{article}

% ============================================================================
% ENCODAGE ET LANGUE
% ============================================================================
\usepackage[utf8]{inputenc}
\usepackage[T1]{fontenc}
\usepackage[french]{babel}

% ============================================================================
% MATHÉMATIQUES
% ============================================================================
\usepackage{amsmath,amssymb,amsfonts,mathtools}

% Commandes mathématiques personnalisées (conventions du cours)
\newcommand{\vx}{\mathbf{x}}          % vecteur x
\newcommand{\vw}{\mathbf{w}}          % vecteur w (poids)
\newcommand{\vy}{\mathbf{y}}          % vecteur y
\newcommand{\mX}{\mathbf{X}}          % matrice X
\newcommand{\mW}{\mathbf{W}}          % matrice W
\newcommand{\yhat}{\hat{y}}           % prédiction
\newcommand{\pscal}[2]{\langle #1, #2 \rangle} % produit scalaire
\newcommand{\gradJ}{\nabla J(\boldsymbol{\theta})} % gradient de J

% ============================================================================
% MISE EN PAGE
% ============================================================================
\usepackage[margin=2.5cm]{geometry}
\usepackage{setspace}
\setstretch{1.15}                      % interligne légèrement aéré

% ============================================================================
% COULEURS
% ============================================================================
\usepackage[dvipsnames,svgnames,x11names]{xcolor}

% Palette principale du projet
\definecolor{ipssiBleu}{HTML}{3B82F6}       % Définitions
\definecolor{ipssiVert}{HTML}{22C55E}       % À retenir
\definecolor{ipssiOrange}{HTML}{F97316}     % Attention / Piège
\definecolor{ipssiViolet}{HTML}{A855F7}     % Intuition
\definecolor{ipssiCyan}{HTML}{06B6D4}       % Exemples
\definecolor{ipssiGrisClair}{HTML}{F3F4F6}  % Fond code
\definecolor{ipssiGrisFonce}{HTML}{1F2937}  % Texte code
\definecolor{ipssiFondPage}{HTML}{FAFAFA}   % Fond léger page de garde
\definecolor{ipssiRouge}{HTML}{E74C3C}      % Classe 0 (salaire ≤ 45k€)
\definecolor{ipssiVertData}{HTML}{2ECC71}   % Classe 1 (salaire > 45k€)
\definecolor{ipssiBleuDroite}{HTML}{3498DB} % Droite de régression

% ============================================================================
% ENCADRÉS PÉDAGOGIQUES (tcolorbox)
% ============================================================================
\usepackage[most]{tcolorbox}

% Style de base partagé
\tcbset{
    ipssibase/.style={
        enhanced,
        breakable,
        left=4mm,
        right=4mm,
        top=3mm,
        bottom=3mm,
        boxrule=0pt,
        frame hidden,
        borderline west={3pt}{0pt}{#1},
        colback=#1!5,
        colframe=#1,
        fonttitle=\bfseries\sffamily,
        coltitle=#1!80!black,
        attach boxed title to top left={yshift=-2mm, xshift=4mm},
        boxed title style={
            colback=#1!15,
            colframe=#1,
            boxrule=0.5pt,
            arc=2pt,
        },
        arc=2pt,
        outer arc=2pt,
    }
}

% Icônes TikZ pour les encadrés (compatible pdflatex, pas d'emoji Unicode)
\newcommand{\icoDefinition}{%
    \tikz[baseline=-0.3em]{\draw[ipssiBleu, thick] (0,0) -- (0.3,0) -- (0.3,0.35) -- (0,0.35) -- cycle;
    \draw[ipssiBleu, thick] (0.05,0.1) -- (0.25,0.1); \draw[ipssiBleu, thick] (0.05,0.2) -- (0.2,0.2);}%
}
\newcommand{\icoRetenir}{%
    \tikz[baseline=-0.2em]{\draw[ipssiVert, thick, fill=ipssiVert!20] (0.15,0) -- (0.3,0.3) -- (0,0.3) -- cycle;}%
}
\newcommand{\icoAttention}{%
    \tikz[baseline=-0.2em]{\draw[ipssiOrange, thick] (0.15,0.35) -- (0,0) -- (0.3,0) -- cycle;
    \node[ipssiOrange, font=\tiny\bfseries] at (0.15,0.12) {!};}%
}
\newcommand{\icoIntuition}{%
    \tikz[baseline=-0.2em]{\draw[ipssiViolet, thick, fill=ipssiViolet!15]
    (0.15,0.35) -- (0.08,0.2) -- (0.08,0.15) -- (0.22,0.15) -- (0.22,0.2) -- cycle;
    \draw[ipssiViolet, thick] (0.11,0.12) -- (0.11,0.05); \draw[ipssiViolet, thick] (0.19,0.12) -- (0.19,0.05);}%
}
\newcommand{\icoExemple}{%
    \tikz[baseline=-0.2em]{\draw[ipssiCyan, thick] (0,0) rectangle (0.3,0.3);
    \draw[ipssiCyan, thick] (0,0.15) -- (0.3,0.15); \draw[ipssiCyan, thick] (0.15,0) -- (0.15,0.3);}%
}
\newcommand{\icoPython}{%
    \tikz[baseline=-0.2em]{\node[font=\small\bfseries\ttfamily, ipssiVert!70!black] at (0,0) {\textgreater\textgreater};}%
}

% Définition (bleu)
\newtcolorbox{definitionbox}[1][D\'{e}finition]{
    ipssibase=ipssiBleu,
    title={\icoDefinition\; #1},
}

% À retenir (vert)
\newtcolorbox{retenirbox}[1][\`{A} retenir]{
    ipssibase=ipssiVert,
    title={\icoRetenir\; #1},
}

% Attention / Piège (orange)
\newtcolorbox{attentionbox}[1][Attention]{
    ipssibase=ipssiOrange,
    title={\icoAttention\; #1},
}

% Intuition (violet)
\newtcolorbox{intuitionbox}[1][Intuition]{
    ipssibase=ipssiViolet,
    title={\icoIntuition\; #1},
}

% Exemple (cyan)
\newtcolorbox{exemplebox}[1][Exemple]{
    ipssibase=ipssiCyan,
    title={\icoExemple\; #1},
}

% Code Python (style terminal)
\newtcolorbox{pythonbox}[1][Code Python]{
    enhanced,
    breakable,
    left=4mm, right=4mm, top=3mm, bottom=3mm,
    boxrule=0.5pt,
    colback=ipssiGrisClair,
    colframe=ipssiGrisFonce!40,
    fonttitle=\bfseries\ttfamily\small,
    coltitle=ipssiGrisFonce,
    title={\icoPython\; #1},
    attach boxed title to top left={yshift=-2mm, xshift=4mm},
    boxed title style={colback=ipssiGrisClair, colframe=ipssiGrisFonce!40, boxrule=0.5pt, arc=2pt},
    arc=2pt, outer arc=2pt,
}

% ============================================================================
% CODE PYTHON (listings)
% ============================================================================
\usepackage{listings}

\lstdefinestyle{python}{
    language=Python,
    basicstyle=\ttfamily\small,
    keywordstyle=\color{ipssiBleu}\bfseries,
    stringstyle=\color{ipssiVert!70!black},
    commentstyle=\color{gray}\itshape,
    numberstyle=\tiny\color{gray},
    numbers=left,
    numbersep=8pt,
    backgroundcolor=\color{ipssiGrisClair},
    frame=none,
    rulecolor=\color{ipssiGrisFonce!20},
    breaklines=true,
    breakatwhitespace=true,
    tabsize=4,
    showstringspaces=false,
    captionpos=b,
    aboveskip=8pt,
    belowskip=8pt,
    xleftmargin=10pt,
    framexleftmargin=10pt,
    morekeywords={as, with, True, False, None, self, np, plt, import, from},
    literate=
        {é}{{\'e}}1 {è}{{\`e}}1 {ê}{{\^e}}1 {ë}{{\"e}}1
        {à}{{\`a}}1 {â}{{\^a}}1
        {ù}{{\`u}}1 {û}{{\^u}}1
        {ô}{{\^o}}1
        {î}{{\^i}}1 {ï}{{\"i}}1
        {ç}{{\c{c}}}1
        {É}{{\'E}}1 {È}{{\`E}}1
        {→}{{$\rightarrow$}}1
        {≈}{{$\approx$}}1
        {≤}{{$\leq$}}1
        {≥}{{$\geq$}}1
        {★}{{$\bigstar$}}1
        {ŷ}{{$\hat{y}$}}1
        {·}{{$\cdot$}}1,
}

\lstset{style=python}

% Style pour le code inline
\newcommand{\pyinline}[1]{\lstinline[style=python,basicstyle=\ttfamily\small]{#1}}

% ============================================================================
% GRAPHIQUES
% ============================================================================
\usepackage{tikz}
\usepackage{pgfplots}
\pgfplotsset{compat=1.18}
\usetikzlibrary{arrows.meta, calc, positioning, shapes.geometric, decorations.pathreplacing}

% ============================================================================
% TABLEAUX
% ============================================================================
\usepackage{booktabs}
\usepackage{tabularx}
\usepackage{multirow}

% ============================================================================
% LISTES
% ============================================================================
\usepackage{enumitem}
\setlist[itemize]{itemsep=2pt, parsep=0pt, topsep=4pt}
\setlist[enumerate]{itemsep=2pt, parsep=0pt, topsep=4pt}

% ============================================================================
% HYPERLIENS
% ============================================================================
\usepackage[
    colorlinks=true,
    linkcolor=ipssiBleu!80!black,
    urlcolor=ipssiBleu!80!black,
    citecolor=ipssiVert!80!black,
    bookmarks=true,
    bookmarksopen=true,
    pdfauthor={Mohamed EL AFRIT},
    pdftitle={Mathématiques appliquées à l'IA — IPSSI},
    pdfsubject={Formation Bachelor 2 — IPSSI 2025-2026},
]{hyperref}

% ============================================================================
% IMAGES
% ============================================================================
\usepackage{graphicx}

% ============================================================================
% VARIABLES PARAMÉTRABLES (jour et thème)
% ============================================================================
\newcommand{\leJour}{X}
\newcommand{\leTheme}{Thème}
\newcommand{\jourNumero}[1]{\renewcommand{\leJour}{#1}}
\newcommand{\jourTheme}[1]{\renewcommand{\leTheme}{#1}}

% ============================================================================
% EN-TÊTE ET PIED DE PAGE (fancyhdr)
% ============================================================================
\usepackage{fancyhdr}

\pagestyle{fancy}
\fancyhf{}

% En-tête
\fancyhead[L]{\small\sffamily\textcolor{gray}{IPSSI — Maths appliquées à l'IA}}
\fancyhead[R]{\small\sffamily\textcolor{gray}{Jour \leJour\ — \leTheme}}
\renewcommand{\headrulewidth}{0.4pt}

% Pied de page
\fancyfoot[L]{\footnotesize\textcolor{gray}{Mohamed EL AFRIT\enspace\textbullet\enspace\url{www.mohamedelafrit.com}\enspace\textbullet\enspace CC BY-NC-SA 4.0}}
\fancyfoot[C]{\footnotesize\textcolor{gray}{\thepage}}
\fancyfoot[R]{\footnotesize\textcolor{gray}{2025–2026}}
\renewcommand{\footrulewidth}{0.2pt}

% Style pour la page de garde (pas d'en-tête)
\fancypagestyle{pagegarde}{
    \fancyhf{}
    \renewcommand{\headrulewidth}{0pt}
    \fancyfoot[C]{\footnotesize\textcolor{gray}{\thepage}}
    \renewcommand{\footrulewidth}{0pt}
}

% ============================================================================
% PAGE DE GARDE — \pagegarde{titre}{jour}{theme}
% ============================================================================
\newcommand{\pagegarde}[3]{%
    \thispagestyle{pagegarde}
    \begin{center}
        \vspace*{1.5cm}

        % Logo / Bandeau institution
        {\Large\sffamily\bfseries\textcolor{ipssiBleu}{IPSSI}}\par
        \vspace{2pt}
        {\normalsize\sffamily 2\textsuperscript{e} année Bachelor Informatique}\par

        \vspace{0.8cm}
        {\color{ipssiBleu}\rule{0.6\textwidth}{1.5pt}}\par
        \vspace{0.8cm}

        % Titre du module
        {\LARGE\sffamily\bfseries Mathématiques appliquées\\[4pt] à l'Intelligence Artificielle}\par

        \vspace{1cm}

        % Titre du document
        \begin{tcolorbox}[
            enhanced,
            colback=ipssiBleu!5,
            colframe=ipssiBleu!50,
            boxrule=1pt,
            arc=4pt,
            width=0.85\textwidth,
            halign=center,
            top=10pt, bottom=10pt,
        ]
            {\huge\sffamily\bfseries\textcolor{ipssiBleu!80!black}{#1}}\par
            \vspace{6pt}
            {\Large\sffamily Jour #2 — #3}
        \end{tcolorbox}

        \vspace{1.5cm}

        % Informations auteur
        {\large\sffamily
            \textbf{Auteur :} Mohamed EL AFRIT\\[4pt]
            \url{www.mohamedelafrit.com}\\[8pt]
            \textbf{Année académique :} 2025–2026
        }\par

        \vspace{1.5cm}

        % Licence Creative Commons
        {\color{ipssiBleu}\rule{0.4\textwidth}{0.5pt}}\par
        \vspace{8pt}
        {\small\sffamily
            \textcopyright\ 2025 Mohamed EL AFRIT — IPSSI\\[2pt]
            Ce contenu est distribué sous licence\\[2pt]
            \textbf{Creative Commons BY-NC-SA 4.0}\\[2pt]
            \url{https://creativecommons.org/licenses/by-nc-sa/4.0/deed.fr}
        }\par
    \end{center}
    \newpage
}

% ============================================================================
% NIVEAUX D'EXERCICES
% ============================================================================

% Étoile compatible pdflatex
\newcommand{\starfull}[1]{\textcolor{#1}{$\bigstar$}}
\newcommand{\starempty}{\textcolor{gray!40}{$\bigstar$}}

% Fondamental (1 étoile)
\newcommand{\exofondamental}{%
    \starfull{ipssiVert}\starempty\starempty\enspace%
    {\small\sffamily\textcolor{ipssiVert}{\textbf{Fondamental}}}%
}

% Intermédiaire (2 étoiles)
\newcommand{\exointermediaire}{%
    \starfull{ipssiOrange}\starfull{ipssiOrange}\starempty\enspace%
    {\small\sffamily\textcolor{ipssiOrange}{\textbf{Interm\'{e}diaire}}}%
}

% Avancé (3 étoiles)
\newcommand{\exoavance}{%
    \starfull{ipssiRouge}\starfull{ipssiRouge}\starfull{ipssiRouge}\enspace%
    {\small\sffamily\textcolor{ipssiRouge}{\textbf{Avanc\'{e}}}}%
}

% Commande d'exercice numéroté
\newcounter{exocount}[section]
\newcommand{\exercice}[1]{%
    \stepcounter{exocount}%
    \vspace{8pt}%
    \noindent\textbf{\sffamily Exercice \theexocount}\enspace — \enspace #1\par
    \vspace{4pt}%
}

% ============================================================================
% COMMANDES UTILITAIRES
% ============================================================================

% Séparateur visuel entre sections
\newcommand{\separateur}{%
    \vspace{6pt}
    \begin{center}
        {\color{ipssiBleu!30}\rule{0.3\textwidth}{0.5pt}}
    \end{center}
    \vspace{6pt}
}

% Mot-clé en gras coloré
\newcommand{\motcle}[1]{\textbf{\textcolor{ipssiBleu!80!black}{#1}}}

% Formule encadrée importante
\newcommand{\formule}[1]{%
    \begin{center}
        \fcolorbox{ipssiBleu!50}{ipssiBleu!5}{%
            \quad$\displaystyle #1$\quad%
        }
    \end{center}
}

% ============================================================================
% FIN DU PRÉAMBULE
% ============================================================================
          % Importer ce préambule
%   \jourNumero{1}                  % Définir le numéro du jour (1, 2, 3 ou 4)
%   \jourTheme{Données et représentation mathématique}
%   \begin{document}
%   \pagegarde{Fiche de synthèse}{1}{Données et représentation mathématique}
%   ...
%   \end{document}
% ============================================================================

% --- Classe de document ---
\documentclass[11pt,a4paper]{article}

% ============================================================================
% ENCODAGE ET LANGUE
% ============================================================================
\usepackage[utf8]{inputenc}
\usepackage[T1]{fontenc}
\usepackage[french]{babel}

% ============================================================================
% MATHÉMATIQUES
% ============================================================================
\usepackage{amsmath,amssymb,amsfonts,mathtools}

% Commandes mathématiques personnalisées (conventions du cours)
\newcommand{\vx}{\mathbf{x}}          % vecteur x
\newcommand{\vw}{\mathbf{w}}          % vecteur w (poids)
\newcommand{\vy}{\mathbf{y}}          % vecteur y
\newcommand{\mX}{\mathbf{X}}          % matrice X
\newcommand{\mW}{\mathbf{W}}          % matrice W
\newcommand{\yhat}{\hat{y}}           % prédiction
\newcommand{\pscal}[2]{\langle #1, #2 \rangle} % produit scalaire
\newcommand{\gradJ}{\nabla J(\boldsymbol{\theta})} % gradient de J

% ============================================================================
% MISE EN PAGE
% ============================================================================
\usepackage[margin=2.5cm]{geometry}
\usepackage{setspace}
\setstretch{1.15}                      % interligne légèrement aéré

% ============================================================================
% COULEURS
% ============================================================================
\usepackage[dvipsnames,svgnames,x11names]{xcolor}

% Palette principale du projet
\definecolor{ipssiBleu}{HTML}{3B82F6}       % Définitions
\definecolor{ipssiVert}{HTML}{22C55E}       % À retenir
\definecolor{ipssiOrange}{HTML}{F97316}     % Attention / Piège
\definecolor{ipssiViolet}{HTML}{A855F7}     % Intuition
\definecolor{ipssiCyan}{HTML}{06B6D4}       % Exemples
\definecolor{ipssiGrisClair}{HTML}{F3F4F6}  % Fond code
\definecolor{ipssiGrisFonce}{HTML}{1F2937}  % Texte code
\definecolor{ipssiFondPage}{HTML}{FAFAFA}   % Fond léger page de garde
\definecolor{ipssiRouge}{HTML}{E74C3C}      % Classe 0 (salaire ≤ 45k€)
\definecolor{ipssiVertData}{HTML}{2ECC71}   % Classe 1 (salaire > 45k€)
\definecolor{ipssiBleuDroite}{HTML}{3498DB} % Droite de régression

% ============================================================================
% ENCADRÉS PÉDAGOGIQUES (tcolorbox)
% ============================================================================
\usepackage[most]{tcolorbox}

% Style de base partagé
\tcbset{
    ipssibase/.style={
        enhanced,
        breakable,
        left=4mm,
        right=4mm,
        top=3mm,
        bottom=3mm,
        boxrule=0pt,
        frame hidden,
        borderline west={3pt}{0pt}{#1},
        colback=#1!5,
        colframe=#1,
        fonttitle=\bfseries\sffamily,
        coltitle=#1!80!black,
        attach boxed title to top left={yshift=-2mm, xshift=4mm},
        boxed title style={
            colback=#1!15,
            colframe=#1,
            boxrule=0.5pt,
            arc=2pt,
        },
        arc=2pt,
        outer arc=2pt,
    }
}

% Icônes TikZ pour les encadrés (compatible pdflatex, pas d'emoji Unicode)
\newcommand{\icoDefinition}{%
    \tikz[baseline=-0.3em]{\draw[ipssiBleu, thick] (0,0) -- (0.3,0) -- (0.3,0.35) -- (0,0.35) -- cycle;
    \draw[ipssiBleu, thick] (0.05,0.1) -- (0.25,0.1); \draw[ipssiBleu, thick] (0.05,0.2) -- (0.2,0.2);}%
}
\newcommand{\icoRetenir}{%
    \tikz[baseline=-0.2em]{\draw[ipssiVert, thick, fill=ipssiVert!20] (0.15,0) -- (0.3,0.3) -- (0,0.3) -- cycle;}%
}
\newcommand{\icoAttention}{%
    \tikz[baseline=-0.2em]{\draw[ipssiOrange, thick] (0.15,0.35) -- (0,0) -- (0.3,0) -- cycle;
    \node[ipssiOrange, font=\tiny\bfseries] at (0.15,0.12) {!};}%
}
\newcommand{\icoIntuition}{%
    \tikz[baseline=-0.2em]{\draw[ipssiViolet, thick, fill=ipssiViolet!15]
    (0.15,0.35) -- (0.08,0.2) -- (0.08,0.15) -- (0.22,0.15) -- (0.22,0.2) -- cycle;
    \draw[ipssiViolet, thick] (0.11,0.12) -- (0.11,0.05); \draw[ipssiViolet, thick] (0.19,0.12) -- (0.19,0.05);}%
}
\newcommand{\icoExemple}{%
    \tikz[baseline=-0.2em]{\draw[ipssiCyan, thick] (0,0) rectangle (0.3,0.3);
    \draw[ipssiCyan, thick] (0,0.15) -- (0.3,0.15); \draw[ipssiCyan, thick] (0.15,0) -- (0.15,0.3);}%
}
\newcommand{\icoPython}{%
    \tikz[baseline=-0.2em]{\node[font=\small\bfseries\ttfamily, ipssiVert!70!black] at (0,0) {\textgreater\textgreater};}%
}

% Définition (bleu)
\newtcolorbox{definitionbox}[1][D\'{e}finition]{
    ipssibase=ipssiBleu,
    title={\icoDefinition\; #1},
}

% À retenir (vert)
\newtcolorbox{retenirbox}[1][\`{A} retenir]{
    ipssibase=ipssiVert,
    title={\icoRetenir\; #1},
}

% Attention / Piège (orange)
\newtcolorbox{attentionbox}[1][Attention]{
    ipssibase=ipssiOrange,
    title={\icoAttention\; #1},
}

% Intuition (violet)
\newtcolorbox{intuitionbox}[1][Intuition]{
    ipssibase=ipssiViolet,
    title={\icoIntuition\; #1},
}

% Exemple (cyan)
\newtcolorbox{exemplebox}[1][Exemple]{
    ipssibase=ipssiCyan,
    title={\icoExemple\; #1},
}

% Code Python (style terminal)
\newtcolorbox{pythonbox}[1][Code Python]{
    enhanced,
    breakable,
    left=4mm, right=4mm, top=3mm, bottom=3mm,
    boxrule=0.5pt,
    colback=ipssiGrisClair,
    colframe=ipssiGrisFonce!40,
    fonttitle=\bfseries\ttfamily\small,
    coltitle=ipssiGrisFonce,
    title={\icoPython\; #1},
    attach boxed title to top left={yshift=-2mm, xshift=4mm},
    boxed title style={colback=ipssiGrisClair, colframe=ipssiGrisFonce!40, boxrule=0.5pt, arc=2pt},
    arc=2pt, outer arc=2pt,
}

% ============================================================================
% CODE PYTHON (listings)
% ============================================================================
\usepackage{listings}

\lstdefinestyle{python}{
    language=Python,
    basicstyle=\ttfamily\small,
    keywordstyle=\color{ipssiBleu}\bfseries,
    stringstyle=\color{ipssiVert!70!black},
    commentstyle=\color{gray}\itshape,
    numberstyle=\tiny\color{gray},
    numbers=left,
    numbersep=8pt,
    backgroundcolor=\color{ipssiGrisClair},
    frame=none,
    rulecolor=\color{ipssiGrisFonce!20},
    breaklines=true,
    breakatwhitespace=true,
    tabsize=4,
    showstringspaces=false,
    captionpos=b,
    aboveskip=8pt,
    belowskip=8pt,
    xleftmargin=10pt,
    framexleftmargin=10pt,
    morekeywords={as, with, True, False, None, self, np, plt, import, from},
    literate=
        {é}{{\'e}}1 {è}{{\`e}}1 {ê}{{\^e}}1 {ë}{{\"e}}1
        {à}{{\`a}}1 {â}{{\^a}}1
        {ù}{{\`u}}1 {û}{{\^u}}1
        {ô}{{\^o}}1
        {î}{{\^i}}1 {ï}{{\"i}}1
        {ç}{{\c{c}}}1
        {É}{{\'E}}1 {È}{{\`E}}1
        {→}{{$\rightarrow$}}1
        {≈}{{$\approx$}}1
        {≤}{{$\leq$}}1
        {≥}{{$\geq$}}1
        {★}{{$\bigstar$}}1
        {ŷ}{{$\hat{y}$}}1
        {·}{{$\cdot$}}1,
}

\lstset{style=python}

% Style pour le code inline
\newcommand{\pyinline}[1]{\lstinline[style=python,basicstyle=\ttfamily\small]{#1}}

% ============================================================================
% GRAPHIQUES
% ============================================================================
\usepackage{tikz}
\usepackage{pgfplots}
\pgfplotsset{compat=1.18}
\usetikzlibrary{arrows.meta, calc, positioning, shapes.geometric, decorations.pathreplacing}

% ============================================================================
% TABLEAUX
% ============================================================================
\usepackage{booktabs}
\usepackage{tabularx}
\usepackage{multirow}

% ============================================================================
% LISTES
% ============================================================================
\usepackage{enumitem}
\setlist[itemize]{itemsep=2pt, parsep=0pt, topsep=4pt}
\setlist[enumerate]{itemsep=2pt, parsep=0pt, topsep=4pt}

% ============================================================================
% HYPERLIENS
% ============================================================================
\usepackage[
    colorlinks=true,
    linkcolor=ipssiBleu!80!black,
    urlcolor=ipssiBleu!80!black,
    citecolor=ipssiVert!80!black,
    bookmarks=true,
    bookmarksopen=true,
    pdfauthor={Mohamed EL AFRIT},
    pdftitle={Mathématiques appliquées à l'IA — IPSSI},
    pdfsubject={Formation Bachelor 2 — IPSSI 2025-2026},
]{hyperref}

% ============================================================================
% IMAGES
% ============================================================================
\usepackage{graphicx}

% ============================================================================
% VARIABLES PARAMÉTRABLES (jour et thème)
% ============================================================================
\newcommand{\leJour}{X}
\newcommand{\leTheme}{Thème}
\newcommand{\jourNumero}[1]{\renewcommand{\leJour}{#1}}
\newcommand{\jourTheme}[1]{\renewcommand{\leTheme}{#1}}

% ============================================================================
% EN-TÊTE ET PIED DE PAGE (fancyhdr)
% ============================================================================
\usepackage{fancyhdr}

\pagestyle{fancy}
\fancyhf{}

% En-tête
\fancyhead[L]{\small\sffamily\textcolor{gray}{IPSSI — Maths appliquées à l'IA}}
\fancyhead[R]{\small\sffamily\textcolor{gray}{Jour \leJour\ — \leTheme}}
\renewcommand{\headrulewidth}{0.4pt}

% Pied de page
\fancyfoot[L]{\footnotesize\textcolor{gray}{Mohamed EL AFRIT\enspace\textbullet\enspace\url{www.mohamedelafrit.com}\enspace\textbullet\enspace CC BY-NC-SA 4.0}}
\fancyfoot[C]{\footnotesize\textcolor{gray}{\thepage}}
\fancyfoot[R]{\footnotesize\textcolor{gray}{2025–2026}}
\renewcommand{\footrulewidth}{0.2pt}

% Style pour la page de garde (pas d'en-tête)
\fancypagestyle{pagegarde}{
    \fancyhf{}
    \renewcommand{\headrulewidth}{0pt}
    \fancyfoot[C]{\footnotesize\textcolor{gray}{\thepage}}
    \renewcommand{\footrulewidth}{0pt}
}

% ============================================================================
% PAGE DE GARDE — \pagegarde{titre}{jour}{theme}
% ============================================================================
\newcommand{\pagegarde}[3]{%
    \thispagestyle{pagegarde}
    \begin{center}
        \vspace*{1.5cm}

        % Logo / Bandeau institution
        {\Large\sffamily\bfseries\textcolor{ipssiBleu}{IPSSI}}\par
        \vspace{2pt}
        {\normalsize\sffamily 2\textsuperscript{e} année Bachelor Informatique}\par

        \vspace{0.8cm}
        {\color{ipssiBleu}\rule{0.6\textwidth}{1.5pt}}\par
        \vspace{0.8cm}

        % Titre du module
        {\LARGE\sffamily\bfseries Mathématiques appliquées\\[4pt] à l'Intelligence Artificielle}\par

        \vspace{1cm}

        % Titre du document
        \begin{tcolorbox}[
            enhanced,
            colback=ipssiBleu!5,
            colframe=ipssiBleu!50,
            boxrule=1pt,
            arc=4pt,
            width=0.85\textwidth,
            halign=center,
            top=10pt, bottom=10pt,
        ]
            {\huge\sffamily\bfseries\textcolor{ipssiBleu!80!black}{#1}}\par
            \vspace{6pt}
            {\Large\sffamily Jour #2 — #3}
        \end{tcolorbox}

        \vspace{1.5cm}

        % Informations auteur
        {\large\sffamily
            \textbf{Auteur :} Mohamed EL AFRIT\\[4pt]
            \url{www.mohamedelafrit.com}\\[8pt]
            \textbf{Année académique :} 2025–2026
        }\par

        \vspace{1.5cm}

        % Licence Creative Commons
        {\color{ipssiBleu}\rule{0.4\textwidth}{0.5pt}}\par
        \vspace{8pt}
        {\small\sffamily
            \textcopyright\ 2025 Mohamed EL AFRIT — IPSSI\\[2pt]
            Ce contenu est distribué sous licence\\[2pt]
            \textbf{Creative Commons BY-NC-SA 4.0}\\[2pt]
            \url{https://creativecommons.org/licenses/by-nc-sa/4.0/deed.fr}
        }\par
    \end{center}
    \newpage
}

% ============================================================================
% NIVEAUX D'EXERCICES
% ============================================================================

% Étoile compatible pdflatex
\newcommand{\starfull}[1]{\textcolor{#1}{$\bigstar$}}
\newcommand{\starempty}{\textcolor{gray!40}{$\bigstar$}}

% Fondamental (1 étoile)
\newcommand{\exofondamental}{%
    \starfull{ipssiVert}\starempty\starempty\enspace%
    {\small\sffamily\textcolor{ipssiVert}{\textbf{Fondamental}}}%
}

% Intermédiaire (2 étoiles)
\newcommand{\exointermediaire}{%
    \starfull{ipssiOrange}\starfull{ipssiOrange}\starempty\enspace%
    {\small\sffamily\textcolor{ipssiOrange}{\textbf{Interm\'{e}diaire}}}%
}

% Avancé (3 étoiles)
\newcommand{\exoavance}{%
    \starfull{ipssiRouge}\starfull{ipssiRouge}\starfull{ipssiRouge}\enspace%
    {\small\sffamily\textcolor{ipssiRouge}{\textbf{Avanc\'{e}}}}%
}

% Commande d'exercice numéroté
\newcounter{exocount}[section]
\newcommand{\exercice}[1]{%
    \stepcounter{exocount}%
    \vspace{8pt}%
    \noindent\textbf{\sffamily Exercice \theexocount}\enspace — \enspace #1\par
    \vspace{4pt}%
}

% ============================================================================
% COMMANDES UTILITAIRES
% ============================================================================

% Séparateur visuel entre sections
\newcommand{\separateur}{%
    \vspace{6pt}
    \begin{center}
        {\color{ipssiBleu!30}\rule{0.3\textwidth}{0.5pt}}
    \end{center}
    \vspace{6pt}
}

% Mot-clé en gras coloré
\newcommand{\motcle}[1]{\textbf{\textcolor{ipssiBleu!80!black}{#1}}}

% Formule encadrée importante
\newcommand{\formule}[1]{%
    \begin{center}
        \fcolorbox{ipssiBleu!50}{ipssiBleu!5}{%
            \quad$\displaystyle #1$\quad%
        }
    \end{center}
}

% ============================================================================
% FIN DU PRÉAMBULE
% ============================================================================


\jourNumero{2}
\jourTheme{Optimisation}

\begin{document}

\pagegarde{Fiche d'exercices}{2}{Optimisation}

\setstretch{1.08}
\setlength{\parskip}{3pt}
\setlength{\abovedisplayskip}{6pt}
\setlength{\belowdisplayskip}{6pt}

% ======================================================================
% RAPPELS
% ======================================================================
\begin{tcolorbox}[
    enhanced, colback=ipssiCyan!5, colframe=ipssiCyan!50,
    boxrule=0.8pt, arc=4pt,
    title={\sffamily\bfseries Rappels --- Formules du Jour 2},
    coltitle=white, colbacktitle=ipssiCyan!70,
]
\small
\textbf{Fonction de co\^{u}t :}
$J(\vw) = \frac{1}{n}\sum_{i=1}^{n}(y_i - \yhat_i)^2$
\hfill
\textbf{Mise \`{a} jour :}
$w \leftarrow w - \alpha \cdot \frac{\partial J}{\partial w}$

\medskip
\textbf{D\'{e}riv\'{e}es utiles :}
$(c)' = 0$,\quad
$(ax+b)' = a$,\quad
$(x^2)' = 2x$,\quad
$(ax^2)' = 2ax$,\quad
$((x{-}a)^2)' = 2(x{-}a)$

\medskip
\textbf{Gradients r\'{e}gression} ($\yhat_i = w_1 x_i + w_0$) :
$\frac{\partial J}{\partial w_1} = \frac{-2}{n}\sum x_i(y_i - \yhat_i)$
\hfill
$\frac{\partial J}{\partial w_0} = \frac{-2}{n}\sum (y_i - \yhat_i)$

\medskip
\textbf{Dataset salaire} (30 dev, r\'{e}gression exp\'{e}rience $\to$ salaire) :
$\yhat = 2{,}07x + 31{,}27$, MSE $= 20{,}46$.
\end{tcolorbox}

% ======================================================================
\section*{Exercices crayon-papier --- Niveau fondamental}
% ======================================================================

% ----------------------------------------------------------------------
\exercice{\exofondamental\quad --- Calculer le MSE}

Un mod\`{e}le a produit les pr\'{e}dictions suivantes :

\begin{center}
\renewcommand{\arraystretch}{1.2}
\begin{tabular}{c|c|c}
    \toprule
    $i$ & Salaire r\'{e}el $y_i$ (k\texteuro) & Pr\'{e}diction $\yhat_i$ (k\texteuro) \\
    \midrule
    1 & 30 & 33 \\
    2 & 42 & 40 \\
    3 & 55 & 52 \\
    4 & 38 & 44 \\
    \bottomrule
\end{tabular}
\end{center}

\begin{enumerate}[label=\alph*)]
    \item Calculez l'erreur $e_i = y_i - \yhat_i$ pour chaque observation.
    \item Calculez l'erreur au carr\'{e} $e_i^2$ pour chaque observation.
    \item Calculez le MSE : $J = \frac{1}{n}\sum e_i^2$.
    \item L'observation $i = 4$ a la plus grande erreur.
        Si on corrigeait cette seule pr\'{e}diction pour qu'elle soit parfaite ($\yhat_4 = 38$),
        quel serait le nouveau MSE ? De combien a-t-il diminu\'{e} ?
\end{enumerate}

% ----------------------------------------------------------------------
\exercice{\exofondamental\quad --- Calculer des d\'{e}riv\'{e}es}

Calculez la d\'{e}riv\'{e}e $f'(x)$ des fonctions suivantes.
Pr\'{e}cisez la r\`{e}gle utilis\'{e}e \`{a} chaque fois.

\begin{enumerate}[label=\alph*)]
    \item $f(x) = 3x^2$
    \item $f(x) = x^2 - 6x + 9$
    \item $f(x) = 2x + 5$
    \item $f(x) = (x - 7)^2$
    \item $f(x) = 5(x - 2)^2 + 3$
\end{enumerate}

Pour les fonctions a), b) et d), calculez la valeur de $x$ o\`{u} $f'(x) = 0$.
Que repr\'{e}sente ce point ?

% ----------------------------------------------------------------------
\exercice{\exofondamental\quad --- Une it\'{e}ration de descente de gradient}

On consid\`{e}re la fonction de co\^{u}t $J(w) = (w - 4)^2$.

\begin{enumerate}[label=\alph*)]
    \item Calculez la d\'{e}riv\'{e}e $J'(w)$.

    \item On part de $w_0 = 0$ avec un taux d'apprentissage $\alpha = 0{,}2$.
        Calculez la d\'{e}riv\'{e}e en $w_0$ : $J'(0) = ?$

    \item Appliquez la r\`{e}gle de mise \`{a} jour :
        $w_1 = w_0 - \alpha \cdot J'(w_0)$.
        Quelle est la nouvelle valeur $w_1$ ?

    \item Calculez $J(w_0)$ et $J(w_1)$. Le co\^{u}t a-t-il diminu\'{e} ?

    \item Le minimum th\'{e}orique est $w^* = 4$. \`{A} quel pourcentage du chemin
        \^{e}tes-vous apr\`{e}s une seule it\'{e}ration ?
\end{enumerate}

% ----------------------------------------------------------------------
\exercice{\exofondamental\quad --- Trois it\'{e}rations compl\`{e}tes}

On consid\`{e}re $J(w) = w^2$ (minimum en $w^* = 0$). Taux d'apprentissage $\alpha = 0{,}1$.

\begin{enumerate}[label=\alph*)]
    \item Calculez $J'(w)$.

    \item En partant de $w_0 = 5$, remplissez le tableau suivant :

    \begin{center}
    \renewcommand{\arraystretch}{1.3}
    \begin{tabular}{c|c|c|c|c}
        \toprule
        It\'{e}ration & $w$ & $J'(w)$ & $w_{\text{new}} = w - \alpha J'(w)$ & $J(w)$ \\
        \midrule
        0 & 5{,}00 & \phantom{00}?\phantom{00} & \phantom{00}?\phantom{00} & \phantom{00}?\phantom{00} \\
        1 & ?      & ?  & ?  & ?  \\
        2 & ?      & ?  & ?  & ?  \\
        \bottomrule
    \end{tabular}
    \end{center}

    \item Apr\`{e}s 3 it\'{e}rations, de combien $w$ s'est-il rapproch\'{e} de $w^* = 0$ ?
        Calculez le ratio $w_3 / w_0$.

    \item \textbf{Bonus :} montrez que $w_k = w_0 \cdot (1 - 2\alpha)^k$ pour cette fonction.
        V\'{e}rifiez avec vos r\'{e}sultats.
\end{enumerate}

% ======================================================================
\section*{Exercices crayon-papier --- Niveau interm\'{e}diaire}
% ======================================================================

% ----------------------------------------------------------------------
\exercice{\exointermediaire\quad --- D\'{e}riv\'{e}es partielles de la r\'{e}gression}

On consid\`{e}re 3 observations avec le mod\`{e}le $\yhat_i = w_1 x_i + w_0$ :

\begin{center}
\renewcommand{\arraystretch}{1.2}
\begin{tabular}{c|c|c}
    \toprule
    $i$ & $x_i$ (exp\'{e}rience) & $y_i$ (salaire, k\texteuro) \\
    \midrule
    1 & 2 & 32 \\
    2 & 5 & 44 \\
    3 & 10 & 55 \\
    \bottomrule
\end{tabular}
\end{center}

\begin{enumerate}[label=\alph*)]
    \item Avec $w_1 = 3$ et $w_0 = 25$, calculez les pr\'{e}dictions $\yhat_i$ et les erreurs $e_i = y_i - \yhat_i$.

    \item Calculez la d\'{e}riv\'{e}e partielle par rapport \`{a} $w_1$ :
    \[
        \frac{\partial J}{\partial w_1} = \frac{-2}{n}\sum_{i=1}^{3} x_i \cdot e_i
    \]

    \item Calculez la d\'{e}riv\'{e}e partielle par rapport \`{a} $w_0$ :
    \[
        \frac{\partial J}{\partial w_0} = \frac{-2}{n}\sum_{i=1}^{3} e_i
    \]

    \item Avec $\alpha = 0{,}01$, calculez les nouveaux poids :
        $w_1' = w_1 - \alpha \frac{\partial J}{\partial w_1}$ et
        $w_0' = w_0 - \alpha \frac{\partial J}{\partial w_0}$.

    \item Dans quelle direction chaque poids a-t-il boug\'{e} ? Est-ce coh\'{e}rent ?
\end{enumerate}

% ----------------------------------------------------------------------
\exercice{\exointermediaire\quad --- Diagnostiquer un probl\`{e}me de convergence}

Un \'{e}tudiant a ex\'{e}cut\'{e} une descente de gradient sur $J(w) = (w - 3)^2$ avec $\alpha = 1{,}5$
et obtient les r\'{e}sultats suivants :

\begin{center}
\renewcommand{\arraystretch}{1.2}
\begin{tabular}{c|c|c}
    \toprule
    It\'{e}ration & $w$ & $J(w)$ \\
    \midrule
    0 & 0{,}00  & 9{,}00 \\
    1 & 9{,}00  & 36{,}00 \\
    2 & $-15{,}00$ & 324{,}00 \\
    3 & 51{,}00  & 2\,304 \\
    4 & $-93{,}00$ & 9\,216 \\
    \bottomrule
\end{tabular}
\end{center}

\begin{enumerate}[label=\alph*)]
    \item V\'{e}rifiez le calcul de l'it\'{e}ration 0 $\to$ 1 :
        $w_1 = w_0 - \alpha \cdot J'(w_0)$.

    \item Le co\^{u}t $J$ augmente-t-il ou diminue-t-il ? Quel est le probl\`{e}me ?

    \item Avec quelle valeur de $\alpha$ la convergence serait-elle assur\'{e}e pour
        cette fonction ? \emph{Indice :} pour $J(w) = (w-a)^2$, la convergence
        n\'{e}cessite $0 < \alpha < 1$.

    \item Recalculez les 3 premi\`{e}res it\'{e}rations avec $\alpha = 0{,}1$.
        Le co\^{u}t diminue-t-il cette fois ?
\end{enumerate}

% ----------------------------------------------------------------------
\exercice{\exointermediaire\quad --- Comparer deux taux d'apprentissage}

On minimise $J(w) = (w - 5)^2$ en partant de $w_0 = 0$.
Deux \'{e}tudiants utilisent des $\alpha$ diff\'{e}rents.

\begin{enumerate}[label=\alph*)]
    \item \textbf{\'{E}tudiant A} ($\alpha = 0{,}05$) : calculez $w_1, w_2, w_3$ et les co\^{u}ts associ\'{e}s.

    \item \textbf{\'{E}tudiant B} ($\alpha = 0{,}3$) : calculez $w_1, w_2, w_3$ et les co\^{u}ts associ\'{e}s.

    \item Remplissez un tableau comparatif :

    \begin{center}
    \renewcommand{\arraystretch}{1.2}
    \begin{tabular}{c|c|c|c|c}
        \toprule
        & $w_1$ & $w_2$ & $w_3$ & $J(w_3)$ \\
        \midrule
        $\alpha = 0{,}05$ & ? & ? & ? & ? \\
        $\alpha = 0{,}3$  & ? & ? & ? & ? \\
        \bottomrule
    \end{tabular}
    \end{center}

    \item Quel \'{e}tudiant se rapproche le plus vite de $w^* = 5$ ? Lequel risque-t-il
        de d\'{e}passer le minimum avec un $\alpha$ plus grand ?
\end{enumerate}

% ======================================================================
\section*{Exercices Python guid\'{e}s --- Niveau interm\'{e}diaire}
% ======================================================================

% ----------------------------------------------------------------------
\exercice{\exointermediaire\quad --- Descente de gradient 1D from scratch \hfill\normalfont\small\textsf{[Python guid\'{e}]}}

Compl\'{e}tez le code pour minimiser $J(w) = (w - 3)^2$ par descente de gradient :

\begin{pythonbox}[Squelette --- Descente de gradient 1D]
\begin{lstlisting}
import numpy as np

# Fonction de cout et sa derivee
def J(w):
    return ______  # A COMPLETER : (w - 3) ** 2

def dJ(w):
    return ______  # A COMPLETER : derivee de J

# Parametres
w = 0.0             # point de depart
alpha = ______      # A COMPLETER : taux d'apprentissage (essayer 0.1)
n_iter = 20         # nombre d'iterations

# Boucle de descente de gradient
for i in range(n_iter):
    grad = ______         # A COMPLETER : calculer la derivee
    w = ______            # A COMPLETER : mise a jour de w
    if i < 5 or i == n_iter - 1:
        print(f"Iter {i+1:2d} : w = {w:.4f}, J = {J(w):.4f}")

print(f"\nResultat : w = {w:.4f} (optimal : 3.0)")
print(f"Erreur : |w - 3| = {abs(w - 3):.6f}")
\end{lstlisting}
\end{pythonbox}

% ----------------------------------------------------------------------
\exercice{\exointermediaire\quad --- Visualiser la convergence \hfill\normalfont\small\textsf{[Python guid\'{e}]}}

Compl\'{e}tez le code pour tracer la courbe de co\^{u}t \`{a} chaque it\'{e}ration :

\begin{pythonbox}[Squelette --- Visualisation de la convergence]
\begin{lstlisting}
import numpy as np
import matplotlib.pyplot as plt

def J(w): return (w - 3) ** 2
def dJ(w): return 2 * (w - 3)

w = -2.0
alpha = 0.15
n_iter = 25

# Sauvegarder l'historique
hist_w = [w]
hist_J = [J(w)]

for i in range(n_iter):
    w = w - alpha * dJ(w)
    hist_w.append(______)   # A COMPLETER
    hist_J.append(______)   # A COMPLETER

# --- Graphique 1 : trajectoire sur la courbe ---
fig, (ax1, ax2) = plt.subplots(1, 2, figsize=(13, 5))

w_range = np.linspace(-3, 8, 200)
ax1.plot(w_range, J(w_range), 'b-', lw=2)
ax1.plot(hist_w, hist_J, 'ro-', markersize=5,
         label='Iterations')
ax1.set_xlabel('w'); ax1.set_ylabel('J(w)')
ax1.set_title('Trajectoire sur la courbe')
ax1.legend(); ax1.grid(True, alpha=0.3)

# --- Graphique 2 : J en fonction du numero d'iteration ---
ax2.plot(range(len(hist_J)), ______, 'g-o',  # A COMPLETER
         markersize=4)
ax2.set_xlabel('______')     # A COMPLETER
ax2.set_ylabel('______')     # A COMPLETER
ax2.set_title('Courbe de convergence')
ax2.grid(True, alpha=0.3)

plt.tight_layout(); plt.show()
\end{lstlisting}
\end{pythonbox}

% ----------------------------------------------------------------------
\exercice{\exointermediaire\quad --- Gradient descent sur le dataset salaire \hfill\normalfont\small\textsf{[Python guid\'{e}]}}

Compl\'{e}tez l'impl\'{e}mentation de la descente de gradient pour la r\'{e}gression
lin\'{e}aire (exp\'{e}rience $\to$ salaire) :

\begin{pythonbox}[Squelette --- R\'{e}gression par descente de gradient]
\begin{lstlisting}
import numpy as np
import matplotlib.pyplot as plt

# Dataset
x = np.array([0,1,1,2,2,0,3,1,4,5,5,6,6,7,7,4,
    8,5,9,10,10,11,12,9,13,11,15,17,20,18], dtype=float)
y = np.array([28.0,26.5,35.0,32.0,38.0,25.0,34.5,42.0,
    38.5,44.0,41.0,48.0,43.5,50.0,39.0,47.0,51.0,40.0,
    52.0,55.0,48.5,58.0,60.0,42.0,62.0,54.0,65.0,68.0,
    72.0,58.0])
n = len(x)

# Parametres initiaux
w1, w0 = 0.0, 0.0
alpha = 0.001
n_iter = 5000
hist_mse = []

for it in range(n_iter):
    # 1. Predictions
    y_pred = ______                      # A COMPLETER

    # 2. Erreurs
    errors = ______                      # A COMPLETER

    # 3. Gradients
    dw1 = -2/n * np.sum(______)          # A COMPLETER
    dw0 = -2/n * np.sum(______)          # A COMPLETER

    # 4. Mise a jour
    w1 = w1 - ______                     # A COMPLETER
    w0 = w0 - ______                     # A COMPLETER

    hist_mse.append(np.mean(errors ** 2))

print(f"Resultat : y = {w1:.4f} * x + {w0:.4f}")
print(f"Attendu  : y = 2.0672 * x + 31.2830")
print(f"MSE final = {hist_mse[-1]:.2f}")

# Graphique
fig, (ax1, ax2) = plt.subplots(1, 2, figsize=(13, 5))
ax1.plot(hist_mse); ax1.set_xlabel('Iteration')
ax1.set_ylabel('MSE'); ax1.set_title('Convergence')
ax1.grid(True, alpha=0.3)

ax2.scatter(x, y, c='blue', s=30)
xl = np.linspace(0, 21, 100)
ax2.plot(xl, ______, 'r-', lw=2,     # A COMPLETER
         label=f'GD: y={w1:.2f}x+{w0:.2f}')
ax2.set_xlabel('Experience'); ax2.set_ylabel('Salaire')
ax2.set_title('Regression'); ax2.legend()
ax2.grid(True, alpha=0.3)
plt.tight_layout(); plt.show()
\end{lstlisting}
\end{pythonbox}

% ======================================================================
\section*{Exercices Python autonomes --- Niveau avanc\'{e}}
% ======================================================================

% ----------------------------------------------------------------------
\exercice{\exoavance\quad --- Comparer 4 taux d'apprentissage \hfill\normalfont\small\textsf{[Python autonome]}}

\'{E}crivez un programme Python complet qui :

\begin{enumerate}[label=\alph*)]
    \item D\'{e}finit la fonction $J(w) = (w - 5)^2$ et sa d\'{e}riv\'{e}e.

    \item Ex\'{e}cute la descente de gradient pour 4 valeurs de $\alpha$ :
        $0{,}001$, $0{,}01$, $0{,}1$ et $0{,}9$.
        Point de d\'{e}part : $w_0 = -3$, nombre d'it\'{e}rations : 50.

    \item Trace \textbf{deux graphiques} c\^{o}te \`{a} c\^{o}te :
    \begin{itemize}
        \item \`{A} gauche : les 4 trajectoires sur la courbe $J(w)$ (une couleur par $\alpha$)
        \item \`{A} droite : les 4 courbes de convergence ($J$ vs num\'{e}ro d'it\'{e}ration)
    \end{itemize}

    \item Affiche un tableau r\'{e}capitulatif :
        pour chaque $\alpha$, la valeur finale de $w$ et le MSE final.

    \item R\'{e}pondez en commentaire : quel $\alpha$ est le meilleur ? Pourquoi ?
        Y a-t-il un $\alpha$ qui diverge ?
\end{enumerate}

\begin{attentionbox}[Consignes]
    Utilisez uniquement NumPy et Matplotlib. Commentez en fran\c{c}ais.
    Utilisez une boucle sur les valeurs de $\alpha$ pour \'{e}viter la r\'{e}p\'{e}tition de code.
\end{attentionbox}

% ----------------------------------------------------------------------
\exercice{\exoavance\quad --- Gradient descent vs scikit-learn \hfill\normalfont\small\textsf{[Python autonome]}}

\'{E}crivez un programme qui compare deux m\'{e}thodes pour trouver la droite de
r\'{e}gression (exp\'{e}rience $\to$ salaire) sur le dataset fil rouge complet :

\begin{enumerate}[label=\alph*)]
    \item \textbf{M\'{e}thode 1 --- From scratch :} impl\'{e}mentez la descente de gradient
        compl\`{e}te (d\'{e}part $w_1 = w_0 = 0$, $\alpha = 0{,}001$, 10\,000 it\'{e}rations).
        Enregistrez $w_1$, $w_0$ et le MSE \`{a} chaque it\'{e}ration.

    \item \textbf{M\'{e}thode 2 --- scikit-learn :} utilisez
        \texttt{LinearRegression} de scikit-learn pour trouver les poids optimaux.
        \emph{Rappel :} \texttt{X} doit \^{e}tre de shape \texttt{(n, 1)} pour scikit-learn.

    \item Affichez un tableau comparatif :

    \begin{center}
    \renewcommand{\arraystretch}{1.2}
    \begin{tabular}{l|c|c|c}
        \toprule
        \textbf{M\'{e}thode} & $w_1$ & $w_0$ & MSE \\
        \midrule
        Gradient descent (10k) & ? & ? & ? \\
        scikit-learn & ? & ? & ? \\
        \bottomrule
    \end{tabular}
    \end{center}

    \item Tracez un graphique avec :
    \begin{itemize}
        \item le nuage de points du dataset,
        \item la droite from scratch (rouge),
        \item la droite scikit-learn (verte),
        \item une l\'{e}gende avec le MSE de chaque m\'{e}thode.
    \end{itemize}

    \item Tracez la courbe de convergence du MSE (it\'{e}rations vs MSE)
        avec une ligne horizontale indiquant le MSE de scikit-learn.

    \item R\'{e}pondez en commentaire :
        les deux m\'{e}thodes trouvent-elles le m\^{e}me r\'{e}sultat ?
        Combien d'it\'{e}rations faut-il approximativement pour que la descente
        de gradient atteigne 99\% de la performance de scikit-learn ?
\end{enumerate}

\begin{intuitionbox}[Indice]
    Pour scikit-learn : \texttt{from sklearn.linear\_model import LinearRegression}.
    L'entr\'{e}e \texttt{X} doit avoir le shape \texttt{(30, 1)} :
    utilisez \texttt{x.reshape(-1, 1)}.
\end{intuitionbox}

% ======================================================================
\vfill
\separateur

\begin{center}
    \small\sffamily
    \textbf{R\'{e}partition :}\quad
    4 exercices \starfull{ipssiVert} fondamentaux \quad
    3 exercices \starfull{ipssiOrange}\starfull{ipssiOrange} interm\'{e}diaires \quad
    3 guid\'{e}s Python \starfull{ipssiOrange}\starfull{ipssiOrange} \quad
    2 autonomes \starfull{ipssiRouge}\starfull{ipssiRouge}\starfull{ipssiRouge}
\end{center}

\begin{center}
    {\footnotesize\sffamily\textcolor{gray}{%
        \textcopyright\ 2025 Mohamed EL AFRIT --- IPSSI \enspace\textbullet\enspace
        \url{www.mohamedelafrit.com} \enspace\textbullet\enspace
        CC BY-NC-SA 4.0
    }}
\end{center}

\end{document}
